\begin{lemma}[Step 1, Subcase 1.2 at a PE zero divider]
Let $G$ be a finite loopless multigraph on finite vertex and edge types,
with the elements of the edge type retained as individual edge identities.
Let $e_0$ be a distinguished edge and let $[k]$ be the color set.  Fix an
arbitrary certified input coloring, a terminal tree-sequence $T$, and a
generating coloring $\varphi$.  Let $\mathcal T$ certify that $T$ is an ETT
on exactly these data, with positive rung number $n$, and assume that
$\mathcal T$ has the maximum property.  Assume also that $\varphi$ is
proper and leaves exactly $e_0$ uncolored.

Assume the following complete lower-rung package.  For every arbitrary
certified input coloring, terminal tree-sequence $U$, generating coloring
$\psi$, and positive-rung maximum ETT $\mathcal U$ on those data whose rung
number $m$ is strictly smaller than $n$, all six of the following hold.
Write $U_m$ and $U_{m+1}$ for the last ladder tree and final series tree of
$\mathcal U$, respectively, and write $D_i$, $\psi_i$, $\Theta_i$, and the
chosen vertex at slot $i$ for its exact indexed series data.
\begin{enumerate}
\item $V(U)$ is elementary under $\psi$.
\item For every proper coloring $\sigma$ leaving exactly $e_0$ uncolored
which is $(U_{m+1},D_m,\psi_m)$-stable, and every two colors $a,b$, if at
least one of $a,b$ is missing on $U_{m+1}$ under $\sigma$, then $a,b$ are
$U_{m+1}$-interchangeable under $\sigma$, with isolated vertex components
included.
\item PE support stability descends with the exact indexed control data,
and equal-first-appearance PE stored sets are disjoint.  More precisely,
if a positive slot $i$ is PE with chosen support $v\in V(U)$, the
first-appearance index of $v$ is $j+1$, and a proper coloring $\sigma$
leaving exactly $e_0$ uncolored is $(U_i,D_i,\psi_i)$-stable, then it is
both $(U(v)-v,D_j,\psi_j)$-stable and $(U_j,D_j,\psi_j)$-stable.  Moreover,
if two distinct positive PE slots have distinct chosen supports in $V(U)$
with the same first-appearance index, then the exact stored-color sets at
those two slots are disjoint.
\item If the terminal slot $m$ is PE with chosen support $v$ and recorded
colors $\gamma_m,\delta_m$, then for every proper coloring $\sigma$ leaving
exactly $e_0$ uncolored and $(U_m,D_m,\psi_m)$-stable, the full
$(\gamma_m,\delta_m)$ incidence component through $v$ meets $V(U_m)$
exactly in $v$.
\item Let $\sigma$ be proper, leave exactly $e_0$ uncolored, and be
$(U_m,D_m,\psi_m)$-stable.  If $\delta$ is defective on $V(U_m)$, if
$v\in V(U)$ is an endpoint in $V(U_m)$ of a $\delta$-colored boundary
edge, and if another such endpoint precedes $v$ in the tree order, then
the following precedence conclusion holds.  Whenever $v$ first appears
no later than a positive slot $i$ which is PE or SE and whose chosen
vertex is $v_i$, either $v=v_i$ or $v\prec_U v_i$.
\item Every proper coloring $\sigma$ leaving exactly $e_0$ uncolored and
$(U_m,D_m,\psi_m)$-stable is a $\psi\bmod U_m$ coloring for
$\mathcal U$.
\end{enumerate}

Let $B$ be a hierarchy base for $\mathcal T$ under $\varphi$, and let $H$
be a good hierarchy over $B$.  Put $d=H.\mathit{dividerCount}$, and fix an
integer $q_0>d$.  Assume uniformly that, for every ETT of rung $n$ built
from an arbitrary certified input and arbitrary generating coloring, if it
has maximum property and carries a good hierarchy with fewer than $q_0$
dividers, then its displayed terminal vertex set is elementary under its
generating coloring.  Assume the large-branch bound
$|\overline\varphi(T_n)|\geq2n+11$, where $T_n$ is the last ladder tree of
$\mathcal T$.  Put $A=H.\mathit{starred}(d)$.  In addition to all five
fields of the good hierarchy, assume separately the terminal certificate
\[
 T\text{ is }
 \left(\bigcup_{\eta\in D^H_{n,d+1}}\Gamma^d_\eta\right)^-
 \text{-closed under }\varphi,
\]
and assume that $V(T)$ is elementary under $\varphi$.

Let $\alpha,\beta\in[k]$ be distinct.  Assume that $\alpha$ is missing and
singleton-closed on $T$ under $\varphi$, and that $\beta$ is missing on
$T$.  For $j=1,2$, let $P_j$ be an exact
$(T,\varphi,\{\alpha,\beta\})$-exit path from the outside end $a_j$ to the
tree end $b_j$.  Assume that the two incidence-walk supports are disjoint,
that $a_1\ne a_2$, that $b_1\ne b_2$, and that $b_1\prec_T b_2$.  Retain
an exact edge identity $f_j$ on $P_j$ which is a $\beta$-colored boundary
edge of $T$ incident with $b_j$.  Suppose $b_2\in V(A)$.

Assume that a color $\theta\in[k]$ has already been selected with
$\theta\ne\alpha$ and $\theta\ne\beta$, that $\theta$ is missing and
singleton-closed on $A$ under $\varphi$, and that it is singleton-closed
on $T$ under $\varphi$.

Finally assume the following explicit no-smaller-presentation consequence
of the chosen global minimum.  Let $\sigma$ be any proper coloring leaving
exactly $e_0$ uncolored which is $(T_n,D_n,\varphi)$-stable.  Let $T'$ be
any tree-sequence carrying an exact corresponding same-rung ETT under
$\sigma$ (and hence its own arbitrary certified replacement input), and
assume that this replacement has maximum property.  For every hierarchy
base of the replacement and every good hierarchy $K$ over that base with
fewer than $d$ dividers, suppose that the replacement satisfies
\[
 |\overline\sigma(T'_n)|\geq2n+11
 \quad\hbox{and}\quad
 T'\text{ is }
 \left(\bigcup_{\eta\in D^K_{n,K.\mathit{dividerCount}+1}}
   \Gamma^{K.\mathit{dividerCount}}_\eta\right)^-
 \text{-closed under }\sigma .
\]
Assume additionally that $d=0$ and that the terminal extension kind is
PE.  Then, for any distinct colors $a,b$ with $a$ missing and singleton-closed
on $T'$ under $\sigma$, there do not exist two exact
$(T',\sigma,\{a,b\})$-exit paths having disjoint incidence-walk supports
and distinct outside ends.  Under all these hypotheses, the displayed
Subcase 1.2 data are contradictory.
\end{lemma}
\begin{proof}
Write $n$ for the rung number, $T_n$ for the last ladder tree, $D_i$ for
the exact cumulative stored-color set at slot $i$, and $\varphi_i$ for
the exact coloring at that slot.  Thus $\varphi=\varphi_n$ by the
generating-coloring field of \noderef{ExtendedTashkinovTree}.  We first
make the two exit legs used in every branch completely explicit.

The starred tree $A$ is an ordered initial segment of $T$.  Unfold
\noderef{TashkinovHierarchy}.  If
$d>0$, then \noderef{HierarchyStarredStage} identifies $A$ with the
$d$-th hierarchy stage, and repeated use of the nestedness field of the
hierarchy gives the required initial-segment map to its terminal stage
$T$.  If $d=0$, the same conclusion follows from the core-inside-first
field and then the nestedness and terminal-stage fields of the hierarchy.
All these are entrywise statements by
\noderef{TreeSequenceInitialSegment}.  Unfold
\noderef{TreeVertexOrder}: there are positions $r<s$ in $T$ displaying
$b_1$ and $b_2$.  Since $b_2\in V(A)$ and the entries of $A$ form an
initial block of those of $T$, its occurrence has position $s$ also in
$A$.  The earlier position $r$ is consequently in that same block, and
so $b_1\in V(A)$ by \noderef{TreeSequenceVertexSet}.

Unfold each given \noderef{TreeExitPath}.  The edge-node immediately
preceding the terminal vertex-node $b_j$ represents an exact boundary
edge of $T$: its other adjacent vertex-node on the path is outside $T$,
because $b_j$ is the only vertex-node of the path in $V(T)$.  By
\noderef{TwoColorIncidenceGraph} that edge has color $\alpha$ or
$\beta$.  It cannot have color $\alpha$, since
\noderef{TreeColorSetClosed} says that the exact $\alpha$-colored
boundary set is empty.  Hence this last boundary edge is
$\beta$-colored.  The separately retained identity $f_j$ has the same
two properties by hypothesis; unfolding \noderef{BoundaryColorEdges}
and \noderef{BoundaryEdges} shows in particular that its end other than
$b_j$ lies outside $T$.  Since $b_j\in V(A)\subseteq V(T)$, the same
identity $f_j$ is also a $\beta$-colored boundary edge of $A$.

We next justify which $\alpha,\theta$ components may be exchanged.  If an
$\alpha$-edge node occurs on $P_j$, it is not incident with $T$: otherwise
one of its incident vertex-nodes on this alternating path would be in
$T$, hence would be $b_j$, and the edge would be an $\alpha$-colored
boundary edge, contrary to the preceding closedness.  More generally,
its full \noderef{TwoColorComponent} for $\alpha,\theta$ contains no
vertex of $T$.  For if it did, take a shortest incidence-graph walk from
the external edge to a vertex-node in $T$ and its first vertex-node in
$T$.  The preceding exact edge identity crosses the boundary of $T$ and,
by \noderef{TwoColorIncidenceGraph}, has color $\alpha$ or $\theta$.
This contradicts respectively the assumed singleton closedness of
$\alpha$ or of $\theta$ on $T$.  The identical first-crossing argument
shows that every $\alpha,\theta$ component containing an outside vertex
is wholly outside $T$.

Exchange $\alpha$ and $\theta$ on every full $\alpha,\theta$ component
outside $T$, and call the result $\nu$.  Finiteness makes this a finite
choice of roots, and the preceding paragraph verifies the outside-support
clause of \noderef{OutsideKempeRecoloring}.  By
\noderef{OutsideKempeRecoloringStable}, $\nu$ is proper, leaves exactly
$e_0$ uncolored, fixes every exact edge incident with $T$, preserves the
missing-color set at each vertex of $T$, and is $(T_n,C,\varphi_n)$-stable
for every chosen control set $C$: after unfolding
\noderef{StableColoring}, restrict its incident-edge and vertex clauses
from $T$ to the ordered initial segment $T_n$.

For $j=1,2$, form a walk $Q_j$ with the same ordered vertex- and
edge-nodes as $P_j$.  Each old $\beta$-edge remains $\beta$ because
$\beta\ne\alpha,\theta$.  Every old $\alpha$-edge lies in one of the
selected outside components just proved and becomes $\theta$.  Thus every
successive pair of nodes of $P_j$ is adjacent in
$I_G(\nu;\theta,\beta)$, so $Q_j$ is a walk there; it is the same simple
node list and has the same positive length and support as $P_j$.
Its outside end still misses one of its two colors.  Explicitly, if
$\beta$ was missing at $a_j$, then the definition
\noderef{MissingColors}, together with \noderef{PresentColors}, says
that no incident edge had color $\beta$; none acquires that color in an
$\alpha,\theta$ exchange, so $\beta$ remains missing.  If instead
$\alpha$ was missing, then no incident edge had color $\alpha$.  The
whole $\alpha,\theta$ component through the outside vertex $a_j$ was
selected, so the $\theta$-edges incident with $a_j$ become $\alpha$ and
the new $\theta$-edges are exactly its old $\alpha$-edges; hence $\theta$
is missing there under $\nu$.  We already proved that both tree ends lie
in $A$, while $a_j\notin V(T)$ implies $a_j\notin V(A)$.  Finally, the
only vertex-node of $Q_j$ in $A$ is $b_j$, because the corresponding
claim for $P_j$ excludes every other vertex even from the larger tree
$T$.  This checks every field of
\[
       \operatorname{TreeExitPath}(A,\nu,\theta,\beta,a_j,b_j,Q_j).
       \tag{1}
\]
The supports remain disjoint, $a_1\ne a_2$, and no claim about separation
of the two transformed full components has been made.

Apply \noderef{HierarchyLEMStable} to the just constructed outside
recoloring.  It supplies a corresponding ETT $\mathcal T^\nu$ displayed
on $T$, a hierarchy base $B^\nu$, a hierarchy $H^\nu$, and a good
hierarchy on it.  The corresponding certificate retains the rung, all
ordered ladder trees, exact connecting-edge sets, stored sets, and kinds;
the same-hierarchy conclusion retains every ordered stage, divider,
stored-color set, and Gamma pair.  The result also supplies maximum
property and transports the given terminal minus-closed certificate.
The stability established above identifies the missing-color sets on
$T_n$ under $\nu$ and $\varphi_n$, so the large bound is unchanged.

By the zero-divider hypothesis, $d=0$.  By \noderef{HierarchyStarredStage}, $A$ is the
actual hierarchy core.  We treat the two possible terminal constructions
separately.

The terminal extension is PE by hypothesis.  Use the transported
corresponding ETT $\mathcal T^\nu$ and its transported base $B^\nu$, not
the obsolete pre-recoloring PE data.  Its PE field gives an actual support
vertex $x$ in $T_n$, an exact closure $R_n$ of the prefix through $x$,
and identifies its core with the ordered union $T_n\vee R_n=A$.  The
prefix contains $x$ by the entrywise definition
\noderef{TreeSequencePrefix}.  The augmentation sequence in
\noderef{TashkinovClosure}, after unfolding
\noderef{TashkinovAugmentationSequence} and
\noderef{TashkinovAugmentation}, consists of ordered initial-segment
steps.  Induction through those steps therefore retains the support
prefix in $R_n$, so $x\in V(R_n)$ as well.  Choose a color $c$ missing at $x$ by the same
$k\geq\Delta(G)+1$ degree count as above.  Thus
\noderef{TreeMissingColors} puts $c$ in both
$\overline\nu(T_n)$ and $\overline\nu(R_n)$.

The pre-step elementarity required next is available for these transported
data.  If $n>1$, apply \noderef{CertifiedETTPredecessor} to
$\mathcal T^\nu$ and then lower assertion (i); if $n=1$, apply the
elementarity conclusion of \noderef{TechnicalTheoremBase} to the exact
initial tuple of the replacement series.  Therefore
\noderef{ParallelCommonMissingClosed}, with $\mathcal T^\nu$ and
$B^\nu$, makes every common missing color, and in particular $c$,
singleton-closed on the exact core $A$.

The retained identity $f_2$ is still a $\beta$-colored boundary edge of
$A$: its inside endpoint is $b_2\in V(A)$, its other endpoint is outside
$T$ and hence outside $A$, and the first outside recoloring fixed its
color.  Consequently $c\ne\beta$.  Moreover $\beta$ is present at $x$.
For if it were missing there, then $x\in V(T_n)\cap V(R_n)$ would put
$\beta$ in both tree-missing sets; a second application of
\noderef{ParallelCommonMissingClosed} would make $A$ singleton-$\beta$-
closed, contradicting the exact boundary identity $f_2$.  Unfolding
\noderef{MissingColors} and \noderef{PresentColors} therefore supplies
an exact $\beta$-colored edge incident with $x$.

If $\theta=c$, put $\rho=\nu$; otherwise exchange $\theta,c$ on all full
components outside $A$.  The first outside recoloring fixes incident
edges and missing sets on $A$, so the original $\theta$-missingness and
singleton closedness on $A$ hold under $\nu$; singleton-$c$-closedness was
just proved.  The first-crossing argument again verifies the outside
support required by \noderef{OutsideKempeRecoloring}.  In the nonidentity
case, apply \noderef{OutsideKempeRecoloringStable} first to $T_n$ with
control set $D_n$ and then to $R_n$ with empty control; components outside
$A$ are outside each constituent.  In the identity case use reflexivity
of stability.  The two resulting clauses are exactly
\noderef{CombinedStableColoring} relative to the transported generating
coloring $\nu$.  They also give properness, exact one-uncoloredness, and
preserve the missing sets at every vertex of $A$.

In particular $c$ remains missing at $x$, hence belongs to
$\overline\rho(A)$.  Every exact edge incident with $A$ is fixed, so the
empty $c$-colored boundary under $\nu$ remains empty under $\rho$; thus
$A$ is singleton-$c$-closed under $\rho$.  The exact $\beta$-edge at $x$
is fixed as well, so $\beta$ remains present there.  Repeating the
edge-by-edge and outside-end argument used to construct (1) gives two
exact $(A,\rho,\{c,\beta\})$-exit paths with the old disjoint supports
and distinct outside ends.

The leaf argument above now applies verbatim in
$I_G(\rho;c,\beta)$.  It certifies, via
\noderef{TwoColorPathComponent}, the full component through $x$ and the
full components through the two outside ends as path components; the
first meets $A$ at $x$, and each latter component meets $A$ along its
exit path.  The spanning-path argument chooses one outside component
distinct from the component through $x$, again without asserting that
the two transformed leg components are mutually distinct.  Apply
conclusion (iii) of \noderef{interchange} to $\mathcal T^\nu$, its PE
base $B^\nu$, the just constructed
\noderef{CombinedStableColoring}, and the proved $c$-missing and
singleton-$c$-closed properties on $A$.  It says that $c,\beta$ are
$A$-interchangeable.  The definition \noderef{TreeInterchangeable} then
equates the two distinct full path components meeting $A$, the final
contradiction.  This completes Subcase 1.2.
\end{proof}
